\Abstract
\label{ch:Zusammenfassung}

In Vitruvius legen Konsistenzerhaltungsregeln genau eine Abbildung zwischen Metamodellen fest.
Alternative Abbildungen erfordern Kopien der Regeln.
Diese Arbeit ergänzt die Reactions Language um Feature Annotationen, sodass ein Präprozessor die Regeln jeder Konfiguration aus einem 150\% Modell ableitet.
Ändern sich die Regeln eines persistierten \mbox{V-SUM}, leitet eine vollständige Migration alle übrigen Modelle aus einem dominanten Modell neu ab, eine selektive Migration nur bei betroffenen Elementen.
Für die umljava-reactions umfasst das 150\% Modell 3266 Codezeilen, neun getrennte Kopien dagegen 23791.
Von 72 selektiven Migrationen zwischen verschiedenen Konfigurationen entsprechen 29 einer Ableitung mit einem leeren V-SUM.
Eine selektive Migration betrifft im Median 8,1\,\% der Regeln, aber 66\,\% des UML-Modells.